% ==============================================================================
% SLIDES - SEMANA 01: INTRODUÇÃO À GOVERNANÇA DE TI, GOVERNANÇA VS GESTÃO & AGREGACAO DE VALOR
% ==============================================================================

% ------------------------------------------------------------------------------
% 1. SLIDE DE CAPA
% ------------------------------------------------------------------------------
\senaicover
  {Introdução à Governança de TI}
  {Governança vs. Gestão de TI \& Agregação de Valor para o Negócio}
  {Unidade Curricular: Governança de TI}
  {Prof. André Souza}

% ------------------------------------------------------------------------------
% 2. SLIDE DE VISÃO GERAL / AGENDA (GRID DE 4 CARDS)
% ------------------------------------------------------------------------------
\begin{senaiframe}{Visão Geral \& Agenda}{Os 4 Eixos da Aula}

  \begin{columns}[t]
    \begin{column}{0.48\textwidth}
      \begin{tcolorbox}[senaibluecard, title={1. Conceito de Governança}]
        \begin{itemize}
          \item Evolução da TI de suporte para elemento estratégico.
          \item Direitos de decisão e estrutura de responsabilidades.
        \end{itemize}
      \end{tcolorbox}
      
      \vspace{4pt}
      
      \begin{tcolorbox}[senaiambercard, title={3. Agregação de Valor}]
        \begin{itemize}
          \item Alinhamento de investimentos com metas de negócio.
          \item Mitigação de riscos e otimização de recursos.
        \end{itemize}
      \end{tcolorbox}
    \end{column}

    \begin{column}{0.48\textwidth}
      \begin{tcolorbox}[senairedcard, title={2. Governança vs. Gestão}]
        \begin{itemize}
          \item Governança: Avaliar, Dirigir e Monitorar (EDM).
          \item Gestão: Planejar, Construir, Executar (APO/BAI/DSS).
        \end{itemize}
      \end{tcolorbox}
      
      \vspace{4pt}
      
      \begin{tcolorbox}[senaigreencard, title={4. Oficina do PDGTI}]
        \begin{itemize}
          \item Formação dos grupos de consultoria em TI.
          \item Escolha da empresa-alvo e escopo inicial.
        \end{itemize}
      \end{tcolorbox}
    \end{column}
  \end{columns}

\end{senaiframe}

% ------------------------------------------------------------------------------
% 3. GOVERNANÇA VS GESTÃO (COBIT 2019)
% ------------------------------------------------------------------------------
\begin{senaiframe}{Módulo 1 • Conceitos Fundamentais}{Governança de TI vs. Gestão de TI}

  \begin{columns}[t]
    \begin{column}{0.48\textwidth}
      \begin{tcolorbox}[senaicard, title={Governança de TI (EDM)}]
        \begin{itemize}
          \item \textbf{Responsável:} Conselho e Alta Diretoria.
          \item \textbf{Foco:} Avaliar necessidades, Dirigir estratégias e Monitorar o desempenho.
          \item \textbf{Pergunta:} *"Estamos fazendo as coisas certas?"*
        \end{itemize}
      \end{tcolorbox}
    \end{column}

    \begin{column}{0.48\textwidth}
      \begin{tcolorbox}[senaibluecard, title={Gestão de TI (APO / BAI / DSS)}]
        \begin{itemize}
          \item \textbf{Responsável:} CIO e Gerentes Técnicos.
          \item \textbf{Foco:} Planejar, construir, executar e operacionalizar os serviços de TI.
          \item \textbf{Pergunta:} *"Estamos fazendo com eficiência?"*
        \end{itemize}
      \end{tcolorbox}
    \end{column}
  \end{columns}

\end{senaiframe}

% ------------------------------------------------------------------------------
% 4. PILARES DE AGREGAÇÃO DE VALOR
% ------------------------------------------------------------------------------
\begin{senaiframe}{Módulo 2 • Valor para o Negócio}{Os 5 Pilares da Agregação de Valor}

  \begin{columns}[t]
    \begin{column}{0.48\textwidth}
      \begin{tcolorbox}[senaibluecard, title={Estratégia \& Entrega}]
        \begin{itemize}
          \item \textbf{Alinhamento Estratégico:} A TI viabiliza e acelera as metas corporativas.
          \item \textbf{Entrega de Valor:} Projetos entregues no prazo, custo e benefício esperado.
        \end{itemize}
      \end{tcolorbox}
    \end{column}

    \begin{column}{0.48\textwidth}
      \begin{tcolorbox}[senaiambercard, title={Riscos, Recursos \& Métricas}]
        \begin{itemize}
          \item \textbf{Gestão de Riscos:} Proteção de ativos e continuidade dos negócios.
          \item \textbf{Gestão de Recursos:} Otimização de capital humano, software e nuvem.
        \end{itemize}
      \end{tcolorbox}
    \end{column}
  \end{columns}

\end{senaiframe}

% ------------------------------------------------------------------------------
% 5. REFERÊNCIAS BIBLIOGRÁFICAS
% ------------------------------------------------------------------------------
\begin{senaiframe}{Referências Bibliográficas}{Fontes \& Leitura Recomendada}

  \begin{tcolorbox}[senaicard, title={Obras de Referência}]
    \begin{itemize}
      \item \textbf{FERNANDES, Aguinaldo Aragon; ABREU, Vladimir Ferraz.} *Implantando a governança de TI: da estratégia à gestão de processos e serviços.* 4. ed. Rio de Janeiro: Brasport, 2014.
      \item \textbf{WEILL, Peter; ROSS, Jeanne W.} *Governança de Tecnologia da Informação.* Rio de Janeiro: Elsevier, 2004.
      \item \textbf{ISACA.} *COBIT 2019 Framework: Governance and Management Objectives.* Schaumburg: ISACA, 2018.
    \end{itemize}
  \end{tcolorbox}

\end{senaiframe}

% ------------------------------------------------------------------------------
% 6. APLICAÇÃO NO PDGTI
% ------------------------------------------------------------------------------
\begin{senaiframe}{Aplicação no PDGTI}{Orientações para o Projeto da Semana}

  \vspace{0.3cm}
  \begin{tcolorbox}[senaigreencard, title={Atividades Práticas da Semana 1}]
    \begin{itemize}
      \item \textbf{Formação dos Grupos:} Organização em consultorias de TI (individual ou duplas).
      \item \textbf{Escolha da Empresa-Alvo:} Seleção de empresa real ou fictícia bem documentada.
      \item \textbf{Definição do Escopo Inicial:} Identificação do setor, porte, missão e viabilidade de acesso aos dados para o PDGTI.
    \end{itemize}
  \end{tcolorbox}

\end{senaiframe}
